\chapter*{ABSTRAK}
\addcontentsline{toc}{chapter}{ABSTRAK}

\begin{center}

    \textbf{Perancangan dan Implementasi Warehouse Management Controller \& Sistem Lokalisasi untuk Prototipe Smart Warehouse} \\
    \vspace{0.4cm}
    oleh \\
    \vspace{0.4cm}
    Raizan Iqbal Resi\\
    18222068\\
\end{center}


\begin{spacing}{1.0}
Sistem manajemen gudang FMCG konvensional menghadapi keterbatasan kritis dalam
sinkronisasi inventaris secara \textit{real-time} dan penelusuran posisi fisik
barang, sehingga operasional rentan terhadap ketidaksesuaian stok dan hambatan
proses manual. Laporan ini memaparkan perancangan, implementasi, dan evaluasi
\textit{Warehouse Management Controller} (WMC) serta sistem lokalisasi berbasis
AprilTag untuk prototipe \textit{smart warehouse}. WMC dirancang dalam bentuk
tiga \textit{microservice} yang dapat di-\textit{deploy} secara independen,
yaitu \textit{Authentication Service}, \textit{Goods Detail Service}, dan
\textit{Fleet Order Service}. Sistem lokalisasi memanfaatkan kamera untuk
memetakan posisi robot dan barang ke dalam \textit{grid} diskret berukuran
8×4 yang mencakup area lantai seluas 2,4×1,2 m. Pengujian \textit{stress test}
terhadap \textit{endpoint} API WMC berhasil dilalui tanpa kegagalan permintaan.
Verifikasi sistem lokalisasi berhasil mendeteksi posisi sel \textit{grid} dan
orientasi robot maupun barang secara akurat pada seluruh kondisi pengujian,
yang turut berkontribusi terhadap tercapainya \textit{Task Fulfilment Accuracy}
penuh dalam skenario penyimpanan dan pengambilan barang. Hasil ini menunjukkan
bahwa integrasi arsitektur \textit{microservice} berbasis \textit{event-driven}
dengan sistem lokalisasi visual mampu meningkatkan akurasi pelacakan posisi dan
efisiensi operasional gudang secara signifikan.

Kata kunci: Smart Warehouse, Warehouse Management System, Microservice, Computer Vision, Indoor Tracking.
\end{spacing}

